% General-dimensional base-field descent and quantitative generation.
\section{一般点復元の基礎体降下と定量的生成}
\label{gcd:sec:descent}

本節では可換模型の\emph{幾何学的一般点での存在}を得た後の議論を整理する。
低rotation積の可積分性により、方程式を元の体上に選ぶ。
方程式と同値核の定義体の区別は、第\ref{gdr:subsec:scope}節(2)の通りである。
必要な低次数積の同定と勾配空間の射影的一意性は
第\ref{sec:explicit-jacobian-recovery}節に証明を収める。

\subsection{次数別flat入力から従う一般ファイバーのAS正則性}

\begin{proposition}[一般点のambientに対する追加仮定の除去]
\label{gcd:prop:automatic-as}
$(R,\mathfrak m,k)$をnoetherian局所整域、$K=\operatorname{Frac}R$とする。
$A_R$は非負次数の中央$R$-代数、$(A_R)_0=R$、各$(A_R)_j$は有限自由で、
$A_R/\mathfrak mA_R\simeq k[x_1,\ldots,x_N]$とする。
このとき$A_K$は左右noetherianな二次Koszul AS正則代数であり、
大域次元とGorenstein parameterは$N$、Hilbert級数は$(1-z)^{-N}$である。
さらに$\Omega_R\in Z(A_R)_d$の剰余が非零なら、$\Omega_R$は左右正則で、
商の各次数は$R$-自由である。
\end{proposition}
\begin{proof}
補題\ref{lp:lem:noetherian}の次数別Krull交叉の議論を$A_R$に適用すると、
左右noetherian性を得る。中央集合$R\setminus\{0\}$による局所化で$A_K$もそうなる。
ここでnoetherian性を有限Hilbert級数だけから推論していない。

補題\ref{lp:lem:resolution}を、augmentation加群$R$と、
$A_R^e=A_R\otimes_RA_R^{\mathrm{op}}$上の対角$A_R$に適用する。
多項式環のKoszul分解と対角分解を持ち上げると、項がそれぞれ
$A_R(-j)^{\binom Nj}$、$A_R^e(-j)^{\binom Nj}$である長さ$N$の分解を得る。
持上げの各段で、剰余上の生成元を持ち上げた写像の全射性を各内部次数の
Nakayamaで示す。有限自由$R$-加群への全射は分裂するので核も各次数で自由であり、
最後の核は剰余が零のため零となる。この構成は任意の通常の底変更に可換である。
したがって$A_K$は二次Koszulで、対角分解から左右大域次元は高々$N$となる。

augmentation分解の$A_R$双対を固定内部次数ごとに取る。
その剰余複体は、内部次数$-N$・コホモロジー次数$N$以外で完全である。
有界有限自由複体の単元行列成分を順に消去すれば、他の内部次数の複体は分裂完全となる。
内部次数$-N$では最後の$R$のみが残る。よって左右で
\begin{equation}\label{gcd:eq:as}
 R\underline{\operatorname{Hom}}_{A_R}(R,A_R)\simeq R(N)[-N]
\end{equation}
となり、底変更後の非零な最上Extから大域次元は正確に$N$である。
次数別自由性はHilbert級数を与え、増大度も多項式的である。
最後に各次数の$\Omega_R$乗法は剰余上で単射である。
列数と同じ大きさの単元小行列式を使えば$R$上でも分裂単射となり、
余核も自由となる。中央性により左右で同じ結論を得る。
\end{proof}

\subsection{smooth properなpairから取り出す低次数積}

$k$を標数零の体、$\mathcal C$をsmooth properな$k$-線形dg圏、
$\rho$をdg Morita自己同値とする。対角を$\Delta$、$\rho$の双加群を$P$と書き、
\begin{equation}\label{gcd:eq:low}
 V=H^0R\operatorname{Hom}_{\mathcal C^e}(\Delta,P),\qquad
 U=H^0R\operatorname{Hom}_{\mathcal C^e}(\Delta,P\otimes_{\mathcal C}^{\mathbf L}P)
\end{equation}
と置く。積$\mu:V\otimes V\to U$は$\rho(\beta)\circ\alpha$である。
通常の三角圏上の自然変換と、この導来自然変換を区別する。

\begin{lemma}[低次数積の体拡大との両立]\label{gcd:lem:basechange}
$V,U$は有限次元であり、任意の体拡大$\ell/k$に対して、その形成と積は底変更と可換である。
\end{lemma}
\begin{proof}
smoothnessにより$\Delta$は$\mathcal C^e$上perfectである。
proper性と$P$の可逆性により、$P$と$P\otimes_{\mathcal C}^{\mathbf L}P$は
基礎体上perfectである。有限cellとretractで$\Delta$を表示すると、
式\eqref{gcd:eq:low}のHom複体も$k$-perfectとなる。
同じ有限表示を使えばHomと平坦な体拡大が可換であり、$H^0$も可換である。
双加群のテンソルと射の合成はこの底変更に自然なので、積も保たれる。
\end{proof}

\begin{theorem}[可換cubic模型の方程式は基礎体上に取れる]
\label{gcd:thm:field-model}
$n\ge3$、$N=n+2$とする。上のpairがある代数閉拡大体$\Omega/k$上で
滑らかなcubic $Y/\Omega$のrotation付きKuznetsov pairとMorita同値であると仮定する。
このとき次を満たす$G\in\operatorname{Sym}^3_k V$が存在する。
\begin{enumerate}
\item $X_G=V(G)\subset\mathbb P(V^*)$は$k$上幾何学的に滑らかで、$X_{G,\Omega}\simeq Y$である。
\item ある有限拡大$k'/k$上で$(\mathcal C_{k'},\rho)\simeq(\operatorname{Ku}(X_{G,k'}),\rho_G)$である。
元の同値に三次周期比較を指定していた場合、それも保存できる。
\item pair $(\mathcal C,\rho)$の任意のrotation付き滑らかなcubic模型は、
共通の代数閉拡大体上で射影同型である。
\end{enumerate}
従って方程式の定義体には有限拡大を必要としない。
一意性と同値核の範囲は第\ref{gdr:subsec:scope}節(2)に従う。
\end{theorem}
\begin{proof}
命題\ref{prop:ej-geometric-low-degree}に用いる
Lin--Zhangのsector計算\cite[Theorem~3.5, Remark~3.6]{LZ}を$\Omega$上で適用する。
三次式の非恒等sectorがコホモロジー次数零へ寄与する内部次数は、
$N$が偶数の場合の$N/2$のみである。$N\ge5$では内部次数$1,2$に干渉しない。
したがって、$\Omega$上で$V,U,\mu$はJacobian環の次数$1,2$と積に一致する。
この計算は標数零の代数閉体上の代数的なfactorization計算である。
補題\ref{gcd:lem:basechange}と忠実平坦性から、$k$上でも
\begin{equation}\label{gcd:eq:gradient-data}
 \dim V=N,\quad \dim U=\binom N2,\quad
 \overline\mu:\operatorname{Sym}^2V\twoheadrightarrow U,\quad
 Q:=\ker\overline\mu,\quad\dim Q=N
\end{equation}
となる。積の対称性も降下する。
$Q$には代数閉包上で共通射影零点がない。これは$\Omega$上で勾配空間となることと
$Y$の滑らかさから従う。

$V$の基底を$x_i$、$Q$の基底を$q_a$とする。可積分性の線形解空間
\begin{equation}\label{gcd:eq:integrability}
 \mathcal L_Q=\left\{A\in\operatorname{Mat}_N(k):
 \sum_a A_{ia}\partial_jq_a=\sum_a A_{ja}\partial_iq_a\quad(i<j)\right\}
\end{equation}
は体拡大と可換である。$Y$の方程式の勾配を$q_a$で展開すると、
$\mathcal L_Q\otimes_k\Omega$に可逆行列が存在する。
$k$-基底$A_1,\ldots,A_r$を選べば
$D(z)=\det(\sum_\nu z_\nu A_\nu)$は零多項式でない。
$k$は無限体なので$D(c)\ne0$となる$c\in k^r$が存在する。
実際、各変数の次数は高々$N$なので$\{0,1,\ldots,N\}^r$のいずれかで非零となる。
$A=\sum c_\nu A_\nu$として
\begin{equation}\label{gcd:eq:model}
 G=\frac13\sum_{i,a}x_iA_{ia}q_a
\end{equation}
と置く。可積分性とEuler等式から$\partial_iG=\sum_aA_{ia}q_a$である。
従って$G$の勾配空間は$Q$、その共通射影零点はなく、$X_G$は幾何学的に滑らかである。

$\Omega$上で$G$と$Y$の方程式は同じ勾配空間を持つ。
命題\ref{prop:ej-projective-uniqueness}の対称三重線形形式の証明は
任意の標数零代数閉体で通るので、両者は線形座標変換で一致する。
この証明はsymmetrizer代数の可換性と非零冪零元の不存在を用い、
固有空間ごとの非零係数を三乗根で吸収するものである。
その座標変換はgraded MFの内部shiftと三次周期性を保つ。
従って$\Omega$上で$(\mathcal C,\rho)$と$X_G$のpairの同値を得る。

$X_G$は既に$k$上にあり、そのKuznetsov成分はsmooth properである。
固定された二つのproper dg模型に補題\ref{rog:lem:descent}を適用する。
perfect同値核、rotation比較、指定周期図式を有限生成$k$-部分代数へ降ろし、
単位・余単位のconeを消す非空開の閉点を選べば、有限拡大$k'/k$上の同値を得る。
第三の主張も、任意のrotation付き模型に式\eqref{gcd:eq:gradient-data}を適用し、
同じ勾配空間の射影的一意性を用いれば従う。
\end{proof}

\begin{remark}[存在証明と実計算の違い]\label{gcd:rem:limits}
この定理は幾何学的な可換模型の存在を\emph{仮定して}、その方程式の定義体を改善する。
任意の非可換圏で次元と積分条件を検査するだけでは圏同値は得られない。
また、$V,U,\mu$の有限性の証明は、その非可換入力での具体行列を計算したという意味ではない。
特殊化の範囲は第\ref{gdr:subsec:scope}節(5)を参照する。
\end{remark}

\subsection{積分空間と基礎体上の方程式の曖昧さ}

積分の線形解空間の次元には、Sebastiani--Thom分解による意味がある。
次の命題は補題\ref{ej:lem:symmetrizer-algebra}と
Hwang\cite[Proposition~2.6, Corollary~2.7]{HwangSymmetrizer}のfiber記述を
基礎体上で用いるために書き直したものである。

\begin{proposition}[symmetrizerによる積分空間の記述]
\label{gcd:prop:integrability-symmetrizer}
$k$を標数零の体、$V$を$N$次元$k$-線形空間とし、
$F\in\operatorname{Sym}^3_kV$は幾何学的に滑らかな射影超曲面を定めるとする。
その勾配空間を$Q\subset\operatorname{Sym}^2V$とし、
\[
 \mathcal I_Q=\{G\in\operatorname{Sym}^3V:\partial_iG\in Q\text{ for all }i\}
\]
と置く。ここで$x_i$は$V$の基底であり、この条件は基底に依らない。
$F(x)=T(x,x,x)$を$V^*$上の対称三重線形形式で表し、
補題\ref{ej:lem:symmetrizer-algebra}の代数を$E_F\subset\operatorname{End}_k(V^*)$とする。
このとき
\begin{equation}\label{gcd:eq:integrability-symmetrizer}
 E_F\xrightarrow{\sim}\mathcal I_Q,\qquad
 a\longmapsto F^a,\qquad F^a(x)=T(ax,x,x)
\end{equation}
は$k$-線形同型である。
式\eqref{gcd:eq:integrability}の$\mathcal L_Q$も$\mathcal I_Q$と線形同型で、
\[
 \dim_k\mathcal L_Q=\dim_kE_F
 =\text{$F_{\overline k}$の最細Sebastiani--Thom分解の成分数}
\]
となる。

勾配空間がちょうど$Q$である三次式全体は$E_F^\times$に対応し、全て幾何学的に滑らかである。
この集合に対し、$b\in E_F^\times$による座標変換$x\mapsto bx$と
$k^\times$による方程式の定数倍だけを許すと、その軌道集合は
\begin{equation}\label{gcd:eq:symmetrizer-form-quotient}
 E_F^\times/\bigl(k^\times(E_F^\times)^3\bigr)
\end{equation}
である。これは全ての射影線形座標変換による同型類の分類ではなく、
圏同値核の降下障害を分類するものでもない。
\end{proposition}
\begin{proof}
$\partial_iF$は$Q$の基底なので、$G\in\mathcal I_Q$に対して
$\nabla G=A\nabla F$となる行列$A$が一意に存在する。
対称三重線形形式へ偏極すると、$G$の対称性は$a=A^{\mathsf t}\in E_F$と同値であり、
$G=F^a$となる。逆に$a\in E_F$なら$F^a$は対称三次式で
$\nabla F^a=a^{\mathsf t}\nabla F$を満たす。これで式\eqref{gcd:eq:integrability-symmetrizer}を得る。
任意の基底$q_a$について、$G\mapsto A$を$\nabla G=A(q_a)_a$で定めると、
可積分性と式\eqref{gcd:eq:model}により$\mathcal I_Q\simeq\mathcal L_Q$である。
次元と分解の関係は補題\ref{ej:lem:symmetrizer-algebra}から従う。

勾配が$Q$の基底となる条件は$a$の可逆性である。
この場合、勾配の共通射影零点は$F$と同じく存在しないので$F^a$も幾何学的に滑らかである。
$a,b\in E_F^\times$に対して、symmetrizer条件と可換性から
\[
 F^a(bx)=T(abx,bx,bx)=T(ab^3x,x,x)=F^{ab^3}(x)
\]
である。従って座標変換と定数倍は$a\mapsto c\,ab^3$、$c\in k^\times$として作用し、
その軌道は式\eqref{gcd:eq:symmetrizer-form-quotient}で与えられる。
\end{proof}

\begin{corollary}[幾何学的に直和分解不能な場合]
\label{gcd:cor:indecomposable-equation}
前命題で$F_{\overline k}$が非自明なSebastiani--Thom分解を持たなければ、
$E_F=k$、$\dim_k\mathcal L_Q=1$である。
従って、同じ勾配空間$Q$を持つ$k$上の三次式は全て$F$の$k^\times$倍である。
\end{corollary}
\begin{proof}
補題\ref{ej:lem:symmetrizer-algebra}と
命題\ref{gcd:prop:integrability-symmetrizer}を適用する。
\end{proof}

この系は、固定した勾配空間から積分した方程式の一意性を述べる。
定理\ref{gcd:thm:field-model}の圏同値が$k$上へ降下することまでは意味しない。

\begin{example}[同じ勾配空間を持つ異なる基礎体上のform]
\label{gcd:ex:forms}
$R=\mathbb C[h]_{(h)}$、$k=\mathbb C(h)$、$L=\mathbb C((h))$、$N\ge5$とし、
\begin{equation}\label{gcd:eq:forms}
 F=\sum_{i=1}^Nx_i^3,\qquad
 G=(1+h)x_1^3+\sum_{i=2}^Nx_i^3
\end{equation}
と置く。両者の勾配空間は$\langle x_1^2,\ldots,x_N^2\rangle$であり、
$R$上で滑らかな射影超曲面族を定め、特殊ファイバーは一致する。
二項級数$(1+h)^{1/3}\in\mathbb C[[h]]^\times$で$x_1$を拡大すれば
$F$を$G$へ移せるので、完備底上、従って$L$上では射影同型である。
しかし、$k$上では射影同型でない。
実際、両者のHessian行列式の零集合は座標超平面の和である。
射影線形同型があればこの配置を保つので、その行列は置換行列と対角行列の積となる。
すると$1+h$と$1$の係数比は$k$内の三乗でなければならないが、
$h+1$に沿う離散付値が$1$であることに矛盾する。
従って勾配空間の復元だけから$k$上の射影同型類の一意性を結論できず、
完備底上の同型から$k$上の同型を結論することもできない。
この例では$E_F\simeq k^N$であり、$G$はその単元$(1+h,1,\ldots,1)$に対応する。
式\eqref{gcd:eq:symmetrizer-form-quotient}での非自明性に加え、上のHessianの議論により
全ての射影線形座標変換を許しても両者は同型でないことを確認している。
これは二つの方程式のformの例であり、圏同値核の降下障害が非零であることを
証明する例ではない。
\end{example}

\subsection{既知の生成定理による最短経路と一様な上界}

\begin{proposition}[三shiftの定量的生成：BFKの帰結]
\label{gcd:prop:bound}
$k$は標数零の代数閉体、$Y\subset\mathbb P^{N-1}_k$は滑らかなcubic、$N\ge5$とする。
任意の非零$E\in\operatorname{Ku}(Y)$について
\begin{equation}\label{gcd:eq:bound}
 \operatorname U\bigl(E\oplus\rho_YE\oplus\rho_Y^2E\bigr)
 \le 2N(N+1)-1
\end{equation}
となる。ここで$\operatorname U$は、有限直和・shift・直和因子を段数零として、
圏の全対象を得るためのcone段数の最小値である。
\end{proposition}
\begin{proof}
Orlov比較により$\rho_Y$はgraded singularity categoryの内部shiftとなる。
Ballard--Favero--Katzarkov\cite[Theorem~5.15]{GCDBFK}は、
$N$変数の次数$d$の孤立斉次超曲面について、任意の非零対象の$d$個のshiftの直和が
生成時間$2N\bigl(N(d-2)+1\bigr)-1$以下のstrong generatorとなることを証明している。
$d=3$を代入すればよい。滑らかな射影cubicのアフィン錐の特異点は原点だけである。
\end{proof}
